\section{Manifesto Policy States as the Running Mechanism}
\label{sec:v8-manifesto-mechanism}

The running object is a manifesto scored on one policy dimension. For
concreteness, take the economic-policy scale in the Benoit benchmark. The
oracle \(f^\ast_{\mathrm{econ}}\) is the expert placement a trusted political
measurement procedure assigns to the full platform. A C-Tree state for a span
is useful only if it carries the evidence that can change that placement:
taxation promises, spending commitments, privatization claims, labor-market
policy, welfare language, fiscal constraints, and qualifications about who
benefits or pays.

A local scalar score is too small a state for this task. One leaf may contain
a tax-cut promise, another may contain a health-spending expansion, and a later
section may explain that both are conditional on private delivery or deficit
rules. The economic position of the platform depends on how those pieces fit
together. A parent state that receives only two child scores cannot recover the
cross-span qualification. A parent state that receives two policy-evidence
states can retain the tax position, the welfare position, and the condition
that changes their joint interpretation.

This is the manifesto analogue of a mergeable state. The query is the
expert-facing score, but the object that composes is richer than the score. In
economic policy, the state may say that a party supports tax reduction for
small businesses, protects health spending, and funds the package through
privatization. In environmental policy, it may carry both emissions targets and
exceptions for energy security. In decentralization, it may need to distinguish
administrative devolution from fiscal autonomy, regional veto power, and
country-specific institutional language. The readout happens at the root; the
state must survive every merge before that readout is meaningful.

The three local laws name the failure modes. C1 fails when a leaf summary drops
rubric-relevant evidence in the original span. C2 fails when a produced state
drifts after another summarization pass, such as turning a conditional tax cut
into an unconditional one. C3 fails when two individually adequate child states
are merged into a parent that loses the interaction between them. The C3
failure is the characteristic manifesto problem: policy meaning often sits in
the relation between sections rather than in either section alone.

Decentralization is the stress dimension in the current results because its
evidence geometry is less uniform than economic policy. A party can endorse
local administration while opposing fiscal devolution; it can support regional
identity while keeping central budget authority; it can discuss federal reform
in country-specific institutional terms. A universal summary trained to serve
six dimensions may preserve the common evidence well and still blur these
distinctions. The local diagnostic gap for decentralization marks exactly this
kind of state-sharing failure.

The mechanism also explains why node labels can matter. A full-document expert
target calibrates the root. A leaf or merge label checks whether a smaller
state still carries the same policy evidence for the same dimension. Manifesto
quasi-sentence codings are a natural source for such checks because they
already mark local policy statements. They do not automatically certify a tree;
they become useful when aggregated or sampled in a way that targets the same
oracle-facing policy dimension as the root score.

The rest of the manuscript formalizes this one picture. A C-Tree fixes the
realized leaf partition and merge schedule. The state map \(g\) creates and
merges policy-evidence states. The readout \(f^\ast\), or a learned proxy
\(\widehat f\), scores those states. The local laws ask whether every realized
state can replace the manifesto span it represents for the chosen expert
policy target.
